\documentclass[aspectratio=169,10pt]{beamer}

% Shared Beamer setup for Fall 2026 course slides.
% Clean Cal Poly-inspired theme: no custom class files or uncommon packages.
\mode<presentation>{
  \usetheme{default}
  \usefonttheme{professionalfonts}
}

\definecolor{CalPolyGreen}{HTML}{154734}
\definecolor{CalPolyGold}{HTML}{BD8B13}
\definecolor{CalPolySage}{HTML}{7FA28F}
\definecolor{CalPolyMint}{HTML}{DDF1E7}
\definecolor{CalPolySand}{HTML}{A29F86}
\definecolor{CalPolyBeige}{HTML}{EFEEDE}
\definecolor{DeckInk}{HTML}{1F2933}
\definecolor{DeckMuted}{HTML}{5F6B7A}
\definecolor{DeckSoft}{HTML}{F7F7F5}

\setbeamersize{text margin left=0.55in,text margin right=0.55in}
\setbeamercolor{normal text}{fg=DeckInk,bg=white}
\setbeamercolor{structure}{fg=CalPolyGreen}
\setbeamercolor{title}{fg=white}
\setbeamercolor{frametitle}{fg=CalPolyGreen,bg=white}
\setbeamercolor{block title}{bg=CalPolySage,fg=white}
\setbeamercolor{block body}{bg=CalPolyMint,fg=black}
\setbeamercolor{alerted text}{fg=CalPolyGold}

\setbeamerfont{title}{size=\Huge,series=\bfseries}
\setbeamerfont{frametitle}{size=\Large,series=\bfseries}
\setbeamerfont{block title}{size=\normalsize,series=\bfseries}

\setbeamertemplate{navigation symbols}{}
\setbeamertemplate{itemize item}{\color{CalPolyGold}$\blacktriangleright$}
\setbeamertemplate{itemize subitem}{\color{CalPolyGreen}$\bullet$}
\setbeamertemplate{footline}{
  \hfill{\color{DeckMuted}\scriptsize\insertframenumber/\inserttotalframenumber}
  \hspace{0.18in}\vspace{0.12in}
}
\setbeamertemplate{frametitle}{
  \vspace{0.18in}
  \hspace{0.02in}\insertframetitle\par
  \vspace{0.08in}
  \noindent\makebox[\linewidth][l]{%
    {\color{CalPolyGold}\rule{0.55in}{2pt}}%
    {\color{CalPolyGreen}\rule{\dimexpr\linewidth-0.55in\relax}{2pt}}%
  }
}

\newcommand{\CourseNumber}{Course}
\newcommand{\CourseTitle}{Course Title}
\newcommand{\LectureNumber}{0}
\newcommand{\LectureTitle}{Lecture Title}
\newcommand{\LectureDate}{Date}
\newcommand{\CourseUnit}{Unit}

\newcommand{\coursenumber}[1]{\renewcommand{\CourseNumber}{#1}}
\newcommand{\coursetitle}[1]{\renewcommand{\CourseTitle}{#1}}
\newcommand{\lecturenumber}[1]{\renewcommand{\LectureNumber}{#1}}
\newcommand{\lecturetitle}[1]{\renewcommand{\LectureTitle}{#1}}
\newcommand{\lecturedate}[1]{\renewcommand{\LectureDate}{#1}}
\newcommand{\courseunit}[1]{\renewcommand{\CourseUnit}{#1}}

\author{Kirk Duran}
\institute{California Polytechnic State University, San Luis Obispo}
\date{\LectureDate}

\newcommand{\makedecktitle}{
  \begingroup
  \setbeamercolor{background canvas}{bg=CalPolyGreen}
  \begin{frame}[plain]
    \vfill
    \begin{center}
      {\usebeamerfont{title}\color{white}\LectureTitle\par}
      \vspace{0.18in}
      {\Large\color{white}Lecture \LectureNumber\par}
    \end{center}
    \vfill
    \noindent{\color{white}\rule{\linewidth}{1.2pt}}\par
    \vspace{0.08in}
    \noindent\colorbox{CalPolyGold}{\makebox[0.24\linewidth][c]{\color{white}\Large\CourseNumber}}
    \hspace{0.12in}{\color{white}\Large Kirk Duran}
  \end{frame}
  \endgroup
}

\newenvironment{keyidea}{\begin{block}{Key idea}}{\end{block}}
\newenvironment{activity}{\begin{block}{In class}}{\end{block}}
\newenvironment{cpdefinition}{\begin{block}{Definition}}{\end{block}}
\newenvironment{exampleblockcp}{
  \begingroup
  \setbeamercolor{block title}{bg=CalPolySand,fg=white}
  \setbeamercolor{block body}{bg=CalPolyBeige,fg=black}
  \begin{block}{Example}
}{
  \end{block}
  \endgroup
}

\newcommand{\imageplaceholder}[2][0.3in]{%
  \noindent\makebox[\linewidth][c]{%
    \fcolorbox{CalPolySage}{CalPolyBeige}{%
      \parbox[c][#1][c]{0.86\linewidth}{%
        \centering
        {\color{DeckMuted}\scriptsize Image placeholder: }%
        {\color{DeckInk}\scriptsize\ttfamily\detokenize{#2}}%
      }%
    }%
  }%
}

\newcommand{\sectionframe}[1]{
  \begingroup
  \setbeamercolor{background canvas}{bg=CalPolyGreen}
  \begin{frame}[plain]
    \vfill
    {\color{white}\LARGE\bfseries #1\par}
    \vspace{0.18in}
    {\color{CalPolyGold}\rule{0.22\paperwidth}{3pt}}
    {\color{white}\rule{0.62\paperwidth}{3pt}}
    \vfill
  \end{frame}
  \endgroup
}

\newcommand{\placeholdernote}{
  \begin{block}{Placeholder}
    Replace these bullets with the final lecture material, examples, and in-class activities.
  \end{block}
}

\usepackage[T1]{fontenc}
\usepackage{lmodern}

\pdfstringdefDisableCommands{\def\translate#1{#1}}

\graphicspath{{../assets/lecture-25/}}

% If an image file is not yet available, the slide displays a labeled
% placeholder using the intended filename. Place the matching file in
% ../assets/lecture-25/ to replace the placeholder automatically.
\newcommand{\lectureimage}[2][1.75in]{%
  \IfFileExists{../assets/lecture-25/#2}{%
    \includegraphics[width=\linewidth,height=#1,keepaspectratio]{#2}%
  }{%
    \fbox{%
      \parbox[c][#1][c]{0.94\linewidth}{%
        \centering
        \small\textbf{Image Placeholder}\\[0.08in]
        \scriptsize\texttt{\detokenize{#2}}
      }%
    }%
  }%
}

\setbeamersize{text margin left=0.42in,text margin right=0.42in}

\coursenumber{CS 1001}
\coursetitle{Introduction to Computing}
\lecturenumber{25}
\lecturetitle{Dictionaries and Associative Data}
\lecturedate{Fall 2026}
\courseunit{10. Structured Data}

% Iteration coverage review across Lectures 22--25:
% 22: for/range, bindings, indentation, empty ranges, accumulator mechanics.
% 23: totals/counts/build/filter, while/sentinels, progress and boundary bugs.
% 24: comprehensions, evaluation order, iterable machinery, refactoring.
% 25: break/continue/else, guarded traversal, stable traversal sources,
%     pair unpacking/enumerate, best-so-far state, dictionary accumulation.
% Optional backup: zip, any/all early-return patterns, boundary-case checklist.
% Formal loop-invariant proofs remain in Lecture 32; nested loops are excluded here.
% Python references:
% https://docs.python.org/3/tutorial/controlflow.html
% https://docs.python.org/3/tutorial/datastructures.html#looping-techniques
% https://docs.python.org/3/library/stdtypes.html#dictionary-view-objects
\begin{document}

\makedecktitle

\begin{frame}{Today}
  \begin{itemize}
    \item Explain why sequential position is not always the natural way to organize data.
    \item Model a dictionary as a collection of \textbf{key--value associations}.
    \item Create dictionaries, retrieve values, add new associations, and update existing ones.
    \item Control early exits and skips, preserve progress, and traverse collections safely.
    \item Iterate over keys, values, and pairs; keep the best entry seen so far.
    \item Combine iteration and dictionaries to construct a character-frequency map.
  \end{itemize}

  \begin{keyidea}
    Dictionaries change how data is identified; loop-control statements change how repeated processing proceeds.
  \end{keyidea}
\end{frame}

\begin{frame}{The Boundary for This Lecture}
  \begin{columns}[T,onlytextwidth]
    \begin{column}{0.47\textwidth}
      \begin{block}{Today}
        \begin{itemize}
          \item dictionary literals
          \item keys and values
          \item lookup by key
          \item adding and updating entries
          \item membership
          \item loop control and safe traversal
          \item dictionary iteration and best-so-far state
          \item frequency maps
        \end{itemize}
      \end{block}
    \end{column}
    \begin{column}{0.47\textwidth}
      \begin{alertblock}{Not today}
        \begin{itemize}
          \item hash-function design
          \item collision-resolution algorithms
          \item asymptotic complexity analysis
          \item nested dictionaries
          \item large catalogs of dictionary methods
          \item implementing a hash table from scratch
        \end{itemize}
      \end{alertblock}
    \end{column}
  \end{columns}
\end{frame}

\sectionframe{Part 1: Why Associative Data?}

\begin{frame}{Sequences Organize Values by Position}
  \begin{columns}[T,onlytextwidth]
    \begin{column}{0.55\textwidth}
      \begin{block}{A list}
        \centering
        \texttt{scores = [88, 94, 79, 91]}
      \end{block}

      \begin{itemize}
        \item Each value is identified by an integer index.
        \item Position is useful when order itself matters.
        \item The natural question is: \textbf{which position do I want?}
      \end{itemize}
    \end{column}
    \begin{column}{0.39\textwidth}
      % Intended visual: a row of indexed list slots labeled 0, 1, 2, 3.
      \lectureimage[2.35in]{sequence-position-model.png}
    \end{column}
  \end{columns}
\end{frame}

\begin{frame}{Sometimes Position Is the Wrong Question}
  Suppose each score belongs to a particular student.

  \begin{block}{Parallel information}
    \texttt{names  = ["Ada", "Grace", "Linus", "Edsger"]}\\
    \texttt{scores = [88, 94, 79, 91]}
  \end{block}

  To find Grace's score, we must first determine Grace's position and then use the same position in \texttt{scores}.

  \begin{center}
    \textbf{name} $\longrightarrow$ \textbf{position} $\longrightarrow$ \textbf{score}
  \end{center}

  \begin{keyidea}
    The data relationship we care about is not really ``index 1 has value 94.'' It is ``Grace has score 94.''
  \end{keyidea}
\end{frame}

\begin{frame}{Associative Data Stores the Relationship Directly}
  \begin{columns}[T,onlytextwidth]
    \begin{column}{0.54\textwidth}
      \begin{block}{Dictionary}
        \texttt{scores = \{"Ada": 88, "Grace": 94,}\\
        \hspace{0.78in}\texttt{"Linus": 79, "Edsger": 91\}}
      \end{block}

      Now the lookup relationship is direct:

      \begin{center}
        \textbf{"Grace"} $\longrightarrow$ \textbf{94}
      \end{center}

      The name is not merely stored as another value. It identifies the associated score.
    \end{column}
    \begin{column}{0.40\textwidth}
      % Intended visual: arrows from meaningful keys (names) to score values.
      \lectureimage[2.35in]{key-value-associations.png}
    \end{column}
  \end{columns}
\end{frame}

\begin{frame}{Definition: Dictionary}
  \begin{cpdefinition}
    A \textbf{dictionary} is a mutable mapping from \textbf{keys} to \textbf{values}. Each key identifies one associated value.
  \end{cpdefinition}

  \begin{center}
    \textbf{key} $\longrightarrow$ \textbf{value}
  \end{center}

  \begin{itemize}
    \item A key is used to identify an association.
    \item A value is the data stored for that key.
    \item The complete association is often called a \textbf{key--value pair}.
  \end{itemize}
\end{frame}

\begin{frame}{Anatomy of a Dictionary}
  \begin{block}{Example}
    \centering
    \texttt{inventory = \{"bolts": 120, "washers": 75, "nuts": 90\}}
  \end{block}

  \begin{center}
    \renewcommand{\arraystretch}{1.3}
    \begin{tabular}{l|l}
      \textbf{Key} & \textbf{Associated value} \\\hline
      \texttt{"bolts"} & \texttt{120} \\
      \texttt{"washers"} & \texttt{75} \\
      \texttt{"nuts"} & \texttt{90}
    \end{tabular}
  \end{center}

  \begin{keyidea}
    A dictionary is organized by association, not by numerical position.
  \end{keyidea}
\end{frame}

\begin{frame}{Dictionary Keys Are Unique}
  A dictionary cannot contain two independent entries for the same key.

  \begin{block}{Conceptual state}
    \texttt{"Ada"} $\rightarrow$ \texttt{88}\\
    \texttt{"Grace"} $\rightarrow$ \texttt{94}
  \end{block}

  If we later associate \texttt{"Ada"} with \texttt{95}, the existing association is updated.

  \begin{center}
    \texttt{"Ada"} $\rightarrow$ \texttt{95}
  \end{center}

  \begin{keyidea}
    One key identifies one current value.
  \end{keyidea}
\end{frame}

\begin{frame}{Dictionary Properties}
  For our purposes, a Python dictionary has four important properties.

  \begin{enumerate}
    \item It stores \textbf{key--value associations}.
    \item Each key is \textbf{unique} within the dictionary.
    \item It is \textbf{mutable}: associations can be added or updated.
    \item Values are retrieved using a \textbf{key}, not a positional index.
  \end{enumerate}

  \begin{alertblock}{Mental-model shift}
    Do not ask ``what is element 2?'' Ask ``what value is associated with this key?''
  \end{alertblock}
\end{frame}

\sectionframe{Part 2: Creating, Reading, and Updating Dictionaries}

\begin{frame}{Creating Dictionaries}
  \begin{block}{Dictionary literal}
    \texttt{ages = \{"Ada": 36, "Grace": 85, "Alan": 41\}}
  \end{block}

  \begin{block}{Empty dictionary}
    \texttt{counts = \{\}}
  \end{block}

  \begin{itemize}
    \item A colon separates each key from its value.
    \item Commas separate key--value pairs.
    \item An empty pair of braces creates an empty dictionary.
  \end{itemize}
\end{frame}

\begin{frame}{Retrieving a Value by Key}
  \begin{block}{Dictionary lookup}
    \texttt{ages = \{"Ada": 36, "Grace": 85, "Alan": 41\}}\\[0.08in]
    \texttt{age = ages["Grace"]}
  \end{block}

  The lookup expression evaluates to the associated value:

  \begin{center}
    \texttt{ages["Grace"]} $\longrightarrow$ \texttt{85}
  \end{center}

  \begin{keyidea}
    Square brackets still mean ``retrieve something,'' but the selector is now a key rather than a positional index.
  \end{keyidea}
\end{frame}

\begin{frame}{Membership Tests Keys}
  The \texttt{in} operator asks whether a key is present.

  \begin{block}{Example}
    \texttt{ages = \{"Ada": 36, "Grace": 85\}}\\[0.08in]
    \texttt{"Ada" in ages} $\longrightarrow$ \texttt{True}\\
    \texttt{"Alan" in ages} $\longrightarrow$ \texttt{False}
  \end{block}

  \begin{alertblock}{Important}
    For a dictionary, \texttt{in} tests \textbf{keys}, not values.
  \end{alertblock}
\end{frame}

\begin{frame}{Guarding a Lookup}
  Looking up a key that is not present causes a \texttt{KeyError}.

  \begin{block}{Safe pattern}
    \texttt{if name in ages:}\\
    \hspace{0.20in}\texttt{age = ages[name]}
  \end{block}

  This uses a familiar conditional structure:

  \begin{center}
    \textbf{check whether the key exists} $\longrightarrow$ \textbf{retrieve only when valid}
  \end{center}

  We will use this same idea when constructing a frequency map.
\end{frame}

\begin{frame}{Adding a New Association}
  Assignment through a new key adds a key--value pair.

  \begin{block}{Example}
    \texttt{ages = \{\}}\\
    \texttt{ages["Ada"] = 36}\\
    \texttt{ages["Grace"] = 85}
  \end{block}

  Afterwards:

  \begin{center}
    \texttt{ages} $\rightarrow$ \texttt{\{"Ada": 36, "Grace": 85\}}
  \end{center}

  \begin{keyidea}
    Dictionary assignment binds a value to a key inside the dictionary.
  \end{keyidea}
\end{frame}

\begin{frame}{Updating an Existing Association}
  Assignment through an existing key replaces the old associated value.

  \begin{block}{Example}
    \texttt{scores = \{"Ada": 88, "Grace": 94\}}\\
    \texttt{scores["Ada"] = 95}
  \end{block}

  Afterwards:

  \begin{center}
    \texttt{scores} $\rightarrow$ \texttt{\{"Ada": 95, "Grace": 94\}}
  \end{center}

  No second \texttt{"Ada"} entry is created.
\end{frame}

\begin{frame}{Add or Update Depends on the Key}
  The same syntax has two possible effects.

  \begin{center}
    \renewcommand{\arraystretch}{1.35}
    \begin{tabular}{l|l}
      \textbf{Statement} & \textbf{Effect} \\\hline
      \texttt{d[new\_key] = value} & add a new association \\
      \texttt{d[existing\_key] = value} & update the existing association
    \end{tabular}
  \end{center}

  \begin{keyidea}
    The dictionary remains a mapping with one current value per key.
  \end{keyidea}
\end{frame}

\begin{frame}{Dictionaries Are Mutable}
  \begin{columns}[T,onlytextwidth]
    \begin{column}{0.53\textwidth}
      \begin{block}{Aliasing example}
        \texttt{first = \{"x": 1\}}\\
        \texttt{second = first}\\
        \texttt{second["x"] = 9}
      \end{block}

      Afterwards, both names refer to the same changed dictionary.

      \begin{center}
        \texttt{first["x"]} $\longrightarrow$ \texttt{9}
      \end{center}
    \end{column}
    \begin{column}{0.41\textwidth}
      % Intended visual: first and second both point to one dictionary before and after mutation.
      \lectureimage[2.30in]{dictionary-aliasing-mutation.png}
    \end{column}
  \end{columns}
\end{frame}

\begin{frame}{Length Means Number of Associations}
  \begin{block}{Example}
    \texttt{inventory = \{"bolts": 120, "washers": 75, "nuts": 90\}}
  \end{block}

  Then:

  \begin{center}
    \texttt{len(inventory)} $\longrightarrow$ \texttt{3}
  \end{center}

  The length is the number of key--value pairs, not the sum of the values and not the largest possible key.
\end{frame}

\begin{frame}{What Can Be a Key?}
  Dictionary keys must be suitable for stable identification.

  \begin{itemize}
    \item Strings are common keys: \texttt{"Ada"}, \texttt{"banana"}, \texttt{"CA"}.
    \item Integers can also be keys: \texttt{101}, \texttt{202}, \texttt{404}.
    \item Mutable lists cannot be dictionary keys.
  \end{itemize}

  \begin{alertblock}{For this course}
    We will primarily use immutable primitive values such as strings and integers as keys. The precise requirement is called \textbf{hashability}; implementation details belong in a later data structures course.
  \end{alertblock}
\end{frame}

\sectionframe{Part 3: Controlling Loop Execution}

\begin{frame}{Normal Iteration Runs the Body to Completion}
  In the loops from Lectures 22--24, each iteration normally follows the same control path.

  \begin{center}
    \textbf{bind next value} $\longrightarrow$ \textbf{execute body} $\longrightarrow$ \textbf{request next value}
  \end{center}

  Sometimes the current value changes what should happen next.

  \begin{itemize}
    \item We may want to \textbf{skip the rest of this iteration}.
    \item We may know enough to \textbf{stop the loop entirely}.
    \item We may want special code to run only when the loop \textbf{finishes normally}.
  \end{itemize}

  \begin{keyidea}
    \texttt{continue}, \texttt{break}, and loop \texttt{else} refine the control flow of an existing \texttt{for} loop; they do not create a new kind of iteration.
  \end{keyidea}
\end{frame}

\begin{frame}{\texttt{continue}: Skip the Rest of This Iteration}
  \begin{columns}[T,onlytextwidth]
    \begin{column}{0.55\textwidth}
      \begin{block}{Example}
        \texttt{words = ["cat", "", "dog"]}\\[0.08in]
        \texttt{for word in words:}\\
        \hspace{0.20in}\texttt{if word == "":}\\
        \hspace{0.40in}\texttt{continue}\\
        \hspace{0.20in}\texttt{print(word)}
      \end{block}

      When \texttt{word == ""}, Python skips the remaining statements in the body and requests the next value.
    \end{column}
    \begin{column}{0.39\textwidth}
      % Intended visual: loop path where continue jumps from the condition directly to the next binding.
      \lectureimage[2.30in]{continue-control-flow.png}
    \end{column}
  \end{columns}
\end{frame}

\begin{frame}{Trace \texttt{continue}}
  \begin{center}
    \renewcommand{\arraystretch}{1.25}
    \begin{tabular}{c|c|c|c}
      \textbf{Iteration} & \textbf{word} & \textbf{word == ""} & \textbf{Printed?} \\\hline
      1 & \texttt{"cat"} & False & \texttt{cat} \\
      2 & \texttt{""} & True & no \\
      3 & \texttt{"dog"} & False & \texttt{dog}
    \end{tabular}
  \end{center}

  \begin{keyidea}
    \texttt{continue} ends only the \textbf{current iteration}. The loop itself remains active.
  \end{keyidea}
\end{frame}


\begin{frame}{With \texttt{while}, \texttt{continue} Can Skip the Update}
  \small
  Both examples start with \texttt{values = [0, 5]}.
  \begin{columns}[T,onlytextwidth]
    \begin{column}{0.47\textwidth}
      \begin{alertblock}{Incorrect: stuck at index 0}
        \texttt{i = 0}\\
        \texttt{while i < len(values):}\\
        \hspace{0.18in}\texttt{if values[i] == 0:}\\
        \hspace{0.36in}\texttt{continue}\\
        \hspace{0.18in}\texttt{print(values[i])}\\
        \hspace{0.18in}\texttt{i = i + 1}
      \end{alertblock}
    \end{column}
    \begin{column}{0.47\textwidth}
      \begin{block}{Advance before a possible skip}
        \texttt{i = 0}\\
        \texttt{while i < len(values):}\\
        \hspace{0.18in}\texttt{value = values[i]}\\
        \hspace{0.18in}\texttt{i = i + 1}\\
        \hspace{0.18in}\texttt{if value == 0:}\\
        \hspace{0.36in}\texttt{continue}\\
        \hspace{0.18in}\texttt{print(value)}
      \end{block}
    \end{column}
  \end{columns}
  The corrected loop prints \texttt{5} and stops. The left loop never reaches its update.
  \begin{keyidea}
    \texttt{continue} returns to a \texttt{while} condition; it does not advance your index. A \texttt{for} loop requests its next item automatically.
  \end{keyidea}
\end{frame}

\begin{frame}{Check the Boundary Before Reading the Element}
  \small
  \begin{block}{Walk until a match or the end}
    \texttt{values = [4, 9, 2]}\\
    \texttt{target = 7}\\
    \texttt{i = 0}\\
    \texttt{while i < len(values) and values[i] != target:}\\
    \hspace{0.20in}\texttt{i = i + 1}
  \end{block}
  \begin{itemize}
    \item \texttt{and} checks its left side first. If it is false, Python skips the right side.
    \item At \texttt{i == len(values)}, the boundary test prevents an invalid lookup.
    \item Reversing the two tests can raise \texttt{IndexError}, including on an empty list.
  \end{itemize}
  \begin{activity}
    Final \texttt{i}: target \texttt{7} $\rightarrow$ \texttt{3}; target \texttt{9} $\rightarrow$ \texttt{1}; empty list $\rightarrow$ \texttt{0}. Which results identify a valid matching index?
  \end{activity}
  % A match exists only if i < len(values) after the loop.
\end{frame}

\begin{frame}{\texttt{break}: Terminate the Loop Immediately}
  \begin{columns}[T,onlytextwidth]
    \begin{column}{0.55\textwidth}
      \begin{block}{Example}
        \texttt{values = [4, 9, -2, 7]}\\[0.08in]
        \texttt{for value in values:}\\
        \hspace{0.20in}\texttt{if value < 0:}\\
        \hspace{0.40in}\texttt{break}\\
        \hspace{0.20in}\texttt{print(value)}
      \end{block}

      The loop prints \texttt{4} and \texttt{9}. When \texttt{-2} is encountered, the loop terminates; \texttt{7} is never bound to \texttt{value} by this loop.
    \end{column}
    \begin{column}{0.39\textwidth}
      % Intended visual: loop path where break exits the loop body and continues after the loop.
      \lectureimage[2.30in]{break-control-flow.png}
    \end{column}
  \end{columns}
\end{frame}

\begin{frame}{\texttt{continue}, \texttt{break}, and \texttt{return} Are Different}
  \begin{center}
    \renewcommand{\arraystretch}{1.35}
    \begin{tabular}{l|l}
      \textbf{Statement} & \textbf{Control effect} \\\hline
      \texttt{continue} & stop the current iteration; request the next value \\
      \texttt{break} & terminate the nearest loop \\
      \texttt{return} & terminate the current function call
    \end{tabular}
  \end{center}

  \begin{alertblock}{Tracing question}
    After the statement executes, ask: \textbf{where does control continue?} At the next iteration, after the loop, or back in the caller?
  \end{alertblock}
\end{frame}

\begin{frame}{A \texttt{for} Loop Can Have an \texttt{else} Clause}
  The \texttt{else} clause runs only when the loop reaches its end \textbf{without executing \texttt{break}}.

  \begin{block}{Search pattern}
    \texttt{target = 7}\\
    \texttt{for value in values:}\\
    \hspace{0.20in}\texttt{if value == target:}\\
    \hspace{0.40in}\texttt{print("found")}\\
    \hspace{0.40in}\texttt{break}\\
    \texttt{else:}\\
    \hspace{0.20in}\texttt{print("not found")}
  \end{block}

  \begin{keyidea}
    The \texttt{else} belongs to the \textbf{loop}. It means ``the loop exhausted its input without being terminated by \texttt{break}.''
  \end{keyidea}
\end{frame}

\begin{frame}{Reasoning About Loop \texttt{else}}
  Suppose \texttt{values = [4, 9, 2, 7]}.

  \begin{center}
    \renewcommand{\arraystretch}{1.28}
    \begin{tabular}{l|l|l}
      \textbf{Target} & \textbf{How loop ends} & \textbf{Does \texttt{else} run?} \\\hline
      \texttt{7} & \texttt{break} executes & no \\
      \texttt{5} & input is exhausted normally & yes
    \end{tabular}
  \end{center}

  \begin{itemize}
    \item \texttt{continue} does not suppress the loop \texttt{else}; the loop may still finish normally later.
    \item An empty input also completes normally, so the loop \texttt{else} runs.
    \item A \texttt{while} loop can also have an \texttt{else}: it runs when the condition becomes false without a \texttt{break}.
  \end{itemize}

  \begin{alertblock}{Why this form exists}
    It is useful when the interesting distinction is \textbf{found something and stopped} versus \textbf{examined everything and never found it}.
  \end{alertblock}
\end{frame}

\sectionframe{Part 4: Iterating over Dictionaries}

\begin{frame}{A \texttt{for} Loop over a Dictionary Visits Keys}
  \begin{block}{Example}
    \texttt{scores = \{"Ada": 88, "Grace": 94, "Linus": 79\}}\\[0.08in]
    \texttt{for name in scores:}\\
    \hspace{0.20in}\texttt{print(name)}
  \end{block}

  During iteration, \texttt{name} is repeatedly bound to dictionary keys.

  \begin{keyidea}
    The basic dictionary iteration model is: \textbf{one key per iteration}.
  \end{keyidea}
\end{frame}

\begin{frame}{Use the Key to Retrieve Its Value}
  \begin{block}{Code}
    \texttt{for name in scores:}\\
    \hspace{0.20in}\texttt{score = scores[name]}\\
    \hspace{0.20in}\texttt{print(name, score)}
  \end{block}

  Conceptually, each iteration does two things:

  \begin{enumerate}
    \item bind the next key to \texttt{name},
    \item use that key to retrieve its associated value.
  \end{enumerate}

  \begin{center}
    \textbf{key} $\longrightarrow$ \textbf{dictionary lookup} $\longrightarrow$ \textbf{value}
  \end{center}
\end{frame}

\begin{frame}{Tracing Dictionary Iteration}
  Suppose:

  \begin{block}{Dictionary}
    \texttt{scores = \{"Ada": 88, "Grace": 94, "Linus": 79\}}
  \end{block}

  \begin{center}
    \renewcommand{\arraystretch}{1.25}
    \begin{tabular}{c|c|c}
      \textbf{Iteration} & \textbf{name} & \textbf{scores[name]} \\\hline
      1 & \texttt{"Ada"} & \texttt{88} \\
      2 & \texttt{"Grace"} & \texttt{94} \\
      3 & \texttt{"Linus"} & \texttt{79}
    \end{tabular}
  \end{center}

  The exact loop body determines what computation is performed with each association.
\end{frame}

\begin{frame}{Explicit Key Iteration with \texttt{keys()}}
  Python also provides an explicit view of the keys.

  \begin{block}{Equivalent styles}
    \texttt{for key in data:}\\
    \hspace{0.20in}\texttt{...}\\[0.12in]
    \texttt{for key in data.keys():}\\
    \hspace{0.20in}\texttt{...}
  \end{block}

  For ordinary key iteration, the shorter first form is usually sufficient.
\end{frame}

\begin{frame}{Iterating over Values}
  Use \texttt{values()} when the keys themselves are not needed.

  \begin{block}{Example}
    \texttt{for score in scores.values():}\\
    \hspace{0.20in}\texttt{print(score)}
  \end{block}

  Repeated bindings are values rather than keys:

  \begin{center}
    \texttt{88} $\rightarrow$ \texttt{94} $\rightarrow$ \texttt{79}
  \end{center}
\end{frame}

\begin{frame}{Iterating over Key--Value Pairs}
  Use \texttt{items()} when both parts of each association are needed.

  \begin{block}{Example}
    \texttt{for name, score in scores.items():}\\
    \hspace{0.20in}\texttt{print(name, score)}
  \end{block}

  Each iteration provides one key--value pair.

  \begin{center}
    \texttt{("Ada", 88)} $\rightarrow$ \texttt{("Grace", 94)} $\rightarrow$ \texttt{("Linus", 79)}
  \end{center}
\end{frame}


\begin{frame}{Two Loop Variables Unpack One Supplied Pair}
  \small
  \begin{block}{One pair, two bindings}
    \texttt{name, score = ("Ada", 88)}
  \end{block}
  In \texttt{for name, score in scores.items():}, each supplied pair is unpacked in the same way. It is still \textbf{one loop}.
  \begin{columns}[T,onlytextwidth]
    \begin{column}{0.52\textwidth}
      \begin{block}{Need a position and a value?}
        \texttt{names = ["Ada", "Grace"]}\\
        \texttt{for index, name in enumerate(names):}\\
        \hspace{0.18in}\texttt{print(index, name)}
      \end{block}
    \end{column}
    \begin{column}{0.42\textwidth}
      \begin{block}{Output}
        \texttt{0 Ada}\\
        \texttt{1 Grace}
      \end{block}
    \end{column}
  \end{columns}
  \begin{keyidea}
    \texttt{items()} supplies key--value pairs; \texttt{enumerate()} supplies index--value pairs, starting at index \texttt{0} by default.
  \end{keyidea}
\end{frame}

\begin{frame}{Updating Values Without Changing Membership}
  \small
  \begin{columns}[T,onlytextwidth]
    \begin{column}{0.48\textwidth}
      \begin{block}{Replace list elements by index}
        \texttt{values = [2, 5, 8]}\\
        \texttt{for index, value in enumerate(values):}\\
        \hspace{0.18in}\texttt{values[index] = value * 2}
      \end{block}
      The same list now contains \texttt{[4, 10, 16]}. Its length stays the same.
    \end{column}
    \begin{column}{0.48\textwidth}
      \begin{block}{Replace values for existing keys}
        \texttt{scores = \{"Ada": 88, "Grace": 94\}}\\
        \texttt{for name in scores:}\\
        \hspace{0.18in}\texttt{scores[name] = scores[name] + 1}
      \end{block}
      Result: \texttt{\{"Ada": 89, "Grace": 95\}}. No keys are added or removed.
    \end{column}
  \end{columns}
  \begin{keyidea}
    These statements assign through an index or key. Rebinding the loop variable alone would not replace a stored element. Existing aliases observe these intentional mutations.
  \end{keyidea}
\end{frame}

\begin{frame}{Keep the Collection You Traverse Stable}
  \small
  \begin{columns}[T,onlytextwidth]
    \begin{column}{0.47\textwidth}
      \begin{alertblock}{Removing during list traversal}
        \texttt{values = [0, 0, 1]}\\
        \texttt{for value in values:}\\
        \hspace{0.18in}\texttt{if value == 0:}\\
        \hspace{0.36in}\texttt{values.remove(value)}
      \end{alertblock}
      Result: \texttt{[0, 1]}. The second zero shifts left and is skipped.
    \end{column}
    \begin{column}{0.47\textwidth}
      \begin{block}{Build a separate result}
        \texttt{values = [0, 0, 1]}\\
        \texttt{kept = [value for value in values}\\
        \hspace{0.60in}\texttt{if value != 0]}
      \end{block}
      \texttt{kept} is \texttt{[1]}; \texttt{values} remains \texttt{[0, 0, 1]}.
    \end{column}
  \end{columns}
  \begin{keyidea}
    Do not add or remove dictionary keys while iterating that dictionary; it can raise \texttt{RuntimeError}. Updating values for existing keys is allowed. Read one collection and build another when membership must change.
  \end{keyidea}
\end{frame}

\begin{frame}{Another Accumulator: The Best Entry Seen So Far}
  \small
  \begin{block}{Return the name with the highest score, or None for empty input}
    \texttt{def highest\_name(scores: dict[str, int]):}\\
    \hspace{0.20in}\texttt{best = None}\\
    \hspace{0.20in}\texttt{for name, score in scores.items():}\\
    \hspace{0.40in}\texttt{if best is None or score > scores[best]:}\\
    \hspace{0.60in}\texttt{best = name}\\
    \hspace{0.20in}\texttt{return best}
  \end{block}
  \begin{itemize}
    \item \texttt{best} identifies the highest-scoring entry processed so far.
    \item \texttt{or} skips the lookup while \texttt{best is None}; the first entry establishes a real candidate.
    \item This handles negative scores too. Equal scores keep the first encountered name because the comparison uses \texttt{>}.
  \end{itemize}
  \begin{activity}
    Predict the result for \texttt{\{"Ada": -8, "Grace": -2, "Linus": -2\}} and for \texttt{\{\}}.
  \end{activity}
  % Answers: "Grace" (first highest-scoring entry), and None.
\end{frame}

\begin{frame}{Choose the Iteration Form That Matches the Task}
  \begin{center}
    \renewcommand{\arraystretch}{1.30}
    \begin{tabular}{l|l}
      \textbf{Need} & \textbf{Typical form} \\\hline
      keys only & \texttt{for key in d:} \\
      values only & \texttt{for value in d.values():} \\
      both key and value & \texttt{for key, value in d.items():}
    \end{tabular}
  \end{center}

  \begin{keyidea}
    Choose the view of the dictionary that directly supplies the information the loop body needs.
  \end{keyidea}
\end{frame}

\begin{frame}{Association Is More Important Than Position}
  Modern Python dictionaries preserve insertion order, but their semantic purpose is still association by key.

  \begin{itemize}
    \item \texttt{d[0]} looks up the key \texttt{0}; it does not mean ``first entry.
    \item We test membership with a key.
    \item We design dictionary algorithms around relationships between keys and values.
  \end{itemize}

  \begin{alertblock}{Design principle}
    Do not choose a dictionary because you need ``the third item.'' Choose it because a meaningful key should identify a value.
  \end{alertblock}
\end{frame}

\sectionframe{Part 5: From Counting to a Frequency Map}

\begin{frame}{Review: Counting One Character}
  With one target character, a single counter works well.

  \begin{block}{Count the letter \texttt{"a"}}
    \texttt{count = 0}\\
    \texttt{for char in word:}\\
    \hspace{0.20in}\texttt{if char == "a":}\\
    \hspace{0.40in}\texttt{count = count + 1}
  \end{block}

  The state is one integer because we are answering one counting question.
\end{frame}

\begin{frame}{What If We Want Every Character Count?}
  We could create one variable per possible character:

  \begin{block}{Poorly scaling idea}
    \texttt{count\_a = 0}\\
    \texttt{count\_b = 0}\\
    \texttt{count\_c = 0}\\
    \texttt{...}
  \end{block}

  But the set of characters appearing in the input is data. We should not need a separate variable written in advance for every possible character.

  \begin{keyidea}
    We need a collection in which the \textbf{character itself} can identify its current count.
  \end{keyidea}
\end{frame}

\begin{frame}{A Dictionary Can Store Many Counters}
  Let each character be a key and its observed frequency be the value.

  \begin{center}
    \renewcommand{\arraystretch}{1.25}
    \begin{tabular}{c|c}
      \textbf{Character key} & \textbf{Count value} \\\hline
      \texttt{"b"} & \texttt{1} \\
      \texttt{"a"} & \texttt{3} \\
      \texttt{"n"} & \texttt{2}
    \end{tabular}
  \end{center}

  For the word \texttt{"banana"}:

  \begin{center}
    \texttt{\{"b": 1, "a": 3, "n": 2\}}
  \end{center}
\end{frame}

\begin{frame}{Frequency-Map Algorithm}
  Begin with no observed characters:

  \begin{block}{Initial state}
    \texttt{counts = \{\}}
  \end{block}

  For each character:

  \begin{enumerate}
    \item if the character is already a key, increase its count;
    \item otherwise, create the key with initial count \texttt{1}.
  \end{enumerate}

  \begin{keyidea}
    The dictionary grows only when a character is observed for the first time. Later occurrences update an existing association.
  \end{keyidea}
\end{frame}

\begin{frame}{Trace: Building the Map for \texttt{"banana"}}
  \begin{center}
    \renewcommand{\arraystretch}{1.22}
    \begin{tabular}{c|c|l}
      \textbf{Step} & \textbf{char} & \textbf{counts after iteration} \\\hline
      1 & \texttt{"b"} & \texttt{\{"b": 1\}} \\
      2 & \texttt{"a"} & \texttt{\{"b": 1, "a": 1\}} \\
      3 & \texttt{"n"} & \texttt{\{"b": 1, "a": 1, "n": 1\}} \\
      4 & \texttt{"a"} & \texttt{\{"b": 1, "a": 2, "n": 1\}} \\
      5 & \texttt{"n"} & \texttt{\{"b": 1, "a": 2, "n": 2\}} \\
      6 & \texttt{"a"} & \texttt{\{"b": 1, "a": 3, "n": 2\}}
    \end{tabular}
  \end{center}

  \begin{center}
    \textbf{new character} $\rightarrow$ add key \qquad
    \textbf{seen character} $\rightarrow$ update value
  \end{center}
\end{frame}

\begin{frame}{The Complete Frequency-Map Function}
  \begin{block}{Function}
    \texttt{def frequency\_map(word: str) -> dict[str, int]:}\\
    \hspace{0.20in}\texttt{counts = \{\}}\\
    \hspace{0.20in}\texttt{for char in word:}\\
    \hspace{0.40in}\texttt{if char in counts:}\\
    \hspace{0.60in}\texttt{counts[char] = counts[char] + 1}\\
    \hspace{0.40in}\texttt{else:}\\
    \hspace{0.60in}\texttt{counts[char] = 1}\\
    \hspace{0.20in}\texttt{return counts}
  \end{block}

  \begin{keyidea}
    This function combines iteration, conditionals, dictionary membership, dictionary lookup, and dictionary mutation into one coherent computation.
  \end{keyidea}
\end{frame}

\begin{frame}{Prediction Challenge}
  Without running the code, predict:

  \begin{block}{Call}
    \centering
    \texttt{frequency\_map("letter")}
  \end{block}

  Determine:
  \begin{enumerate}
    \item the order in which characters are encountered,
    \item when a new key is created,
    \item when an existing value is updated,
    \item the final key--value associations.
  \end{enumerate}

  \begin{alertblock}{Do not skip the trace}
    The final dictionary is useful, but the important reasoning is how the state changes after each character.
  \end{alertblock}
\end{frame}


\begin{frame}{Check the Meaning of the State, Including Empty Input}
  \small
  After each iteration of \texttt{frequency\_map}, each stored count describes occurrences in the \textbf{prefix processed so far}.
  \begin{itemize}
    \item Before the loop, no characters have been processed: \texttt{counts = \{\}}.
    \item Processing one character increments exactly its count or starts it at \texttt{1}.
    \item At the end, the counts describe the entire input. The input string is unchanged.
  \end{itemize}
  \begin{center}
    \renewcommand{\arraystretch}{1.2}
    \begin{tabular}{l|l}
      \textbf{Input} & \textbf{Expected result} \\\hline
      \texttt{""} & \texttt{\{\}} \\
      \texttt{"a"} & \texttt{\{"a": 1\}} \\
      \texttt{"aaaa"} & \texttt{\{"a": 4\}} \\
      \texttt{"letter"} & \texttt{\{"l": 1, "e": 2, "t": 2, "r": 1\}}
    \end{tabular}
  \end{center}
  We iterate over the \textbf{word}, so growing the separate \textbf{counts} dictionary is safe.
  % This is informal state reasoning. Formal loop invariants remain in Lecture 32.
\end{frame}

\sectionframe{Part 6: A Preview of the Data Structure Underneath}

\begin{frame}{Why Dictionary Lookup Feels Different from List Search}
  \small
  A hash table uses a key to help locate its stored association.

  \begin{columns}[T,onlytextwidth]
    \begin{column}{0.52\textwidth}
      \begin{itemize}
        \item Lists fundamentally organize data by position.
        \item Dictionaries fundamentally organize data by key.
        \item Python dictionaries are implemented using a \textbf{hash-table} strategy.
        \item The key is processed to help determine where its association is stored.
      \end{itemize}

      \begin{alertblock}{Preview only}
        How hashing works, how collisions are handled, and why the structure is efficient are topics for a later data structures course.
      \end{alertblock}
    \end{column}
    \begin{column}{0.42\textwidth}
      % Intended visual: key -> hash operation -> storage location -> key/value pair.
      \lectureimage[2.15in]{hash-table-preview.png}
    \end{column}
  \end{columns}
\end{frame}

\begin{frame}{Sequence or Dictionary?}
  \begin{center}
    \renewcommand{\arraystretch}{1.32}
    \begin{tabular}{l|l|l}
      \textbf{Question} & \textbf{Sequence} & \textbf{Dictionary} \\\hline
      How is data identified? & position & key \\
      Typical access & \texttt{values[i]} & \texttt{data[key]} \\
      Natural use & ordered data & named associations \\
      Repeated processing & iterate values & iterate keys / values / pairs \\
      Example & daily temperatures & character frequencies
    \end{tabular}
  \end{center}

  \begin{keyidea}
    Choosing a data structure begins with identifying the relationship the program needs to represent.
  \end{keyidea}
\end{frame}

\begin{frame}{Where This Course Stops}
  In CS 1001, you should be able to \textbf{use} a dictionary accurately:

  \begin{itemize}
    \item create key--value associations;
    \item retrieve and update values by key;
    \item test whether a key exists;
    \item iterate over keys, values, and pairs;
    \item use a dictionary to represent naturally associative data.
  \end{itemize}

  \begin{block}{Later in data structures}
    You will ask how a hash table represents these associations internally, how operations are implemented, and how performance changes under different conditions.
  \end{block}
\end{frame}

\begin{frame}{Takeaways}
  \begin{enumerate}
    \item Dictionaries organize data by \textbf{association} rather than numerical position.
    \item Assignment through a key either adds a new association or updates an existing one.
    \item \texttt{continue} skips the remainder of one iteration; \texttt{break} terminates the loop; loop \texttt{else} runs only after normal exhaustion.
    \item Iterating over a dictionary naturally exposes its keys; \texttt{values()} and \texttt{items()} expose other useful views.
    \item Guard lookups, preserve loop progress, and keep traversal sources stable; accumulated state can be counts or a best-so-far candidate.
    \item The same tracing discipline still applies: follow bindings, control flow, mutations, and program state one step at a time.
  \end{enumerate}

  \begin{block}{Next}
    Lecture 26 moves from one structured collection to \textbf{nested data and two-dimensional lists}.
  \end{block}
\end{frame}



\appendix
\sectionframe{Optional Iteration Reference}

\begin{frame}{\texttt{zip}: Traverse Two Sequences Together}
  \small
  \begin{block}{Build associations from paired input}
    \texttt{names = ["Ada", "Grace"]}\\
    \texttt{marks = [88, 94]}\\
    \texttt{scores = \{\}}\\
    \texttt{for name, mark in zip(names, marks):}\\
    \hspace{0.20in}\texttt{scores[name] = mark}
  \end{block}
  Result: \texttt{\{"Ada": 88, "Grace": 94\}}.
  \begin{itemize}
    \item Each iteration receives one pair: \texttt{("Ada", 88)}, then \texttt{("Grace", 94)}.
    \item This visits corresponding positions together; it does not visit every possible combination.
    \item By default, \texttt{zip} stops when the shortest input ends. Check equal lengths first if every value must have a partner.
  \end{itemize}
\end{frame}

\begin{frame}{Any and All Questions Have Different Early Exits}
  \small
  \begin{columns}[T,onlytextwidth]
    \begin{column}{0.48\textwidth}
      \begin{block}{Does any value satisfy the condition?}
        \texttt{def any\_negative(values):}\\
        \hspace{0.16in}\texttt{for value in values:}\\
        \hspace{0.32in}\texttt{if value < 0:}\\
        \hspace{0.48in}\texttt{return True}\\
        \hspace{0.16in}\texttt{return False}
      \end{block}
      One matching value proves \texttt{True}; \texttt{False} requires exhausting the input.
    \end{column}
    \begin{column}{0.48\textwidth}
      \begin{block}{Do all values satisfy the condition?}
        \texttt{def all\_nonnegative(values):}\\
        \hspace{0.16in}\texttt{for value in values:}\\
        \hspace{0.32in}\texttt{if value < 0:}\\
        \hspace{0.48in}\texttt{return False}\\
        \hspace{0.16in}\texttt{return True}
      \end{block}
      One counterexample proves \texttt{False}; \texttt{True} requires exhausting the input.
    \end{column}
  \end{columns}
  \begin{keyidea}
    On \texttt{[]}, these return \texttt{False} and \texttt{True}, respectively. ``No value violates the rule'' does not imply the input contained a value.
  \end{keyidea}
  These are loop patterns; Python also provides the built-ins \texttt{any} and \texttt{all}.
\end{frame}

\begin{frame}{Boundary Cases to Check for a Loop Function}
  \begin{itemize}
    \item \textbf{No iterations:} empty input or an initially false \texttt{while} condition.
    \item \textbf{Exactly one iteration:} does initialization produce the right answer?
    \item \textbf{Conditional processing:} no matches, all matches, first match, and last match.
    \item \textbf{Extremum selection:} negative values and tied best values.
    \item \textbf{Early control transfer:} does each skip preserve progress, and where does each exit continue?
    \item \textbf{Input contract:} is the source left unchanged when a fresh result is promised?
  \end{itemize}
  \begin{keyidea}
    Trace the meaning of the state as well as the statements. A loop can terminate and still compute the wrong result.
  \end{keyidea}
\end{frame}

\end{document}
